% R911: gerado; alterar dados e gerador, não esta prancha.
\clearpage
\refstepcounter{figure}\label{fig:AR06-structure}
\noindent{\large\sffamily\bfseries\color{navy}AR06 · Diagrama de arquitetura}\par\smallskip
\noindent\textbf{Memória metacognitiva, observação e leitura visual}\par
\noindent\textit{Como recuperar experiências sem transformar memória em verdade ou instrução?}\par\medskip
\begin{center}

\tikzset{
 atlasnode/.style={draw=steel,fill=blueBg,rounded corners=2mm,align=center,
  text width=4cm,inner sep=5pt,font=\fontsize{9.5}{11.8}\selectfont,text=textmain},
 atlasroot/.style={atlasnode,fill=navy,text=white,text width=2.8cm,font=\fontsize{10}{12}\selectfont\bfseries},
 atlasdecision/.style={diamond,aspect=2.1,draw=amber,fill=amberBg,align=center,
  text width=3cm,inner sep=2pt,font=\fontsize{9.5}{11.5}\selectfont},
 atlascall/.style={-Stealth,draw=steel,thick},
 atlasevent/.style={-Stealth,draw=violet,dashed,thick},
 atlasdata/.style={-Stealth,draw=teal,densely dotted,thick},
 atlasstructure/.style={draw=slate,thick},
 atlasboundary/.style={-Stealth,draw=amber,dash dot,thick},
 atlasedgelabel/.style={fill=papel,inner sep=1.3pt,font=\fontsize{9}{10}\selectfont}
}

\begin{tikzpicture}
\node[atlasnode] (c0) at (-4.3,0) {\textbf{C1: Preparação contextual}\\Recupera contexto e análise delimitada};
\node[atlasnode] (c1) at (4.3,0) {\textbf{C2: Busca na memória}\\Seleciona registros locais por consulta};
\node[atlasnode] (c2) at (-4.3,-3.0) {\textbf{C3: Reflexão de conclusão}\\Registra episódio sem validar a conclusão};
\node[atlasnode] (c3) at (4.3,-3.0) {\textbf{C4: Persistência de observações}\\Salva observação ou explicita falha};
\node[atlasnode] (c4) at (-4.3,-6.0) {\textbf{C5: Eventos metacognitivos}\\Distribui registros de preparação};
\node[atlasnode] (c5) at (4.3,-6.0) {\textbf{C6: Fonte da rede visual}\\Converte estado local em vista da rede};
\draw[atlascall] (c0.east) to[out=0,in=180] node[atlasedgelabel,pos=.25]{E1} (c1.west);
\draw[atlasevent] (c0.west) -- (-7.0,0) -- node[atlasedgelabel,pos=.25]{E2} (-7.0,-6.0) -- (c4.west);
\draw[atlascall] (c2.east) to[out=0,in=180] node[atlasedgelabel,pos=.25]{E3} (c3.west);
\draw[atlasdata] (c3.north) -- node[atlasedgelabel,pos=.25]{E4} (c1.south);
\draw[atlasdata] (c4.east) to[out=0,in=180] node[atlasedgelabel,pos=.25]{E5} (c5.west);
\end{tikzpicture}
\end{center}
{\small\textbf{Figura \thefigure.} Diagrama de arquitetura --- Memória metacognitiva, observação e leitura visual}\par\smallskip
{\footnotesize\textbf{Chave das relações.} E1: C1 para C2 --- recupera registros (chamada); E2: C1 para C5 --- metacognition.prepared (evento); E3: C3 para C4 --- grava episódio (chamada); E4: C4 para C2 --- registros persistidos (dado); E5: C5 para C6 --- leitura do log JSONL (dado).}\par\smallskip
\noindent\textbf{Como funciona e por quê.} Localize as seis responsabilidades C1--C6. Os identificadores E1--E5 permitem conferir o significado de cada conexão na chave das relações.\par
\noindent\textbf{Justificativa.} Separar componentes e tipos de relação ajuda a localizar onde uma dependência, um evento ou uma fronteira foi declarado.\par
\noindent\textbf{Limite.} Metacognição aqui é mecanismo computacional local e heurístico. Não valida referências, não identifica causa raiz e não mede consciência. A confiança do Trust Engine segue outro trajeto. Uma animação ou snapshot não prova inferência; simulações devem permanecer identificadas.\par\smallskip
\fontecore{mci/metacognitive_loop.py}{MetacognitiveLoop.prepare}
\fontecore{mci/metacognitive_loop.py}{MetacognitiveLoop.complete}
